\chapter{Referencia de métricas de negocio}\label{ch:metrics-catalog}

% Alcance: referencia pedagógica estructurada --- fórmula, desglose de términos, qué mide,
%   por qué importa, referencia, trampa y ejemplo trabajado para las 34 métricas.
%   Seis secciones: Economía unitaria, SaaS e ingresos recurrentes, Embudo y marketing,
%   Plataformas de publicidad, Producto y engagement, Métricas de vanidad.
%   Estilo del libro: prosa fluida bajo encabezados; recuadros example y admonition.
%   No se usa metriccard/mcfield/multicols en este archivo.

Este capítulo es una referencia estructurada, no un argumento. Cada métrica recibe una
definición, una fórmula con cada término explicado, un enunciado en lenguaje llano de lo
que el número mide en realidad, la decisión que informa, una referencia de orden de magnitud
donde exista y la trampa que hace que la métrica mienta. Las definiciones siguen a
\textcite{bendle2020marketing}; las derivaciones viven en los capítulos correspondientes ---
CAC y LTV en \cref{ch:customer-economics}, la mecánica de ROAS/CPA en
\cref{ch:paid-acquisition-platforms}, las métricas de ingresos recurrentes en
\cref{ch:business-models}. Una referencia es un ancla de orden de magnitud, nunca un
objetivo en sí misma.


%% ============================================================
\section{Economía unitaria}

La economía unitaria responde una pregunta antes que cualquier otra: ¿gana dinero una sola
transacción --- un cliente, un pedido, un contrato --- y en qué medida? Si la unidad es
correcta, la escala es una palanca; si es incorrecta, la escala es un acelerador de pérdidas.

\subsection{CAC --- Costo de adquisición de clientes}

\textbf{Fórmula.}
\[
  \mathrm{CAC} = \frac{\text{Sales \& marketing spend}}{\text{New customers acquired}}
\]

\textbf{Qué significa cada término.}
\begin{description}
  \item[Sales \& marketing spend] Todos los costos en un período en que incurre el negocio
    para ganar nuevos clientes --- medios pagados, salarios del equipo de ventas (incluidas
    prestaciones y cargas sociales), suscripciones a herramientas, producción de contenido,
    agencias y gastos generales atribuibles. Completamente cargado; \emph{no} solo gasto en
    publicidad.
  \item[New customers acquired] Clientes pagantes (no prospectos, pruebas ni registros)
    ganados en el mismo período utilizado para la cifra de gasto.
\end{description}

\textbf{Qué mide.} El precio completamente cargado que paga el negocio para convertir un
prospecto en cliente pagante.

\textbf{Por qué importa.} El CAC es una de las dos mitades de la prueba de viabilidad de
adquisición. Comparado con el LTV (\cref{eq:ltv-cac}), establece el límite de cuán
agresivamente puede gastar la empresa antes de que la economía unitaria se vuelva negativa.
Un negocio que no conoce su CAC completamente cargado no puede saber si crea o destruye valor
con cada venta. La derivación completa aparece en \cref{ch:customer-economics}.

\textbf{Referencia.} Juzgarlo únicamente contra el LTV (objetivo LTV:CAC $\ge 3$). Un «CAC
bueno» como cifra aislada no existe --- el mismo \money{600} es excelente para un acuerdo
empresarial y ruinoso para una suscripción de consumo de \money{10}/mes.

\textbf{Trampa.} Excluir los salarios de ventas y \emph{marketing}, las herramientas y los
honorarios de agencias subestima el CAC e infla silenciosamente la razón LTV:CAC --- el
truco más común de embellecimiento de economía unitaria en los \emph{decks} de etapa
temprana.

\begin{example}[CAC para la SaaS de referencia]
Gasto mensual en la SaaS B2B de referencia: medios pagados \money{90000} + salarios
completamente cargados de AE/SDR \money{25000} + herramientas \money{5000} = \money{120000}.
Nuevos clientes pagantes ganados en el mes: 200.
\[
  \mathrm{CAC} = \frac{\money{120000}}{200} = \money{600}.
\]
Excluir salarios y herramientas haría caer la cifra reportada a \money{450} --- una
subestimación del \qty{25}{\percent} que inflaría la razón LTV:CAC de 5.25$\times$ a
7$\times$, posiblemente desencadenando un aumento del presupuesto de adquisición que destruye
valor con la economía unitaria real.
\end{example}


\subsection{LTV --- Customer Lifetime Value}

\textbf{Fórmula.}
\[
  \mathrm{LTV} = \frac{m}{r - g}
\]
donde esta es la forma de perpetuidad creciente (\cref{eq:gordon}) evaluada a nivel de
cliente.

\textbf{Qué significa cada término.}
\begin{description}
  \item[$m$] Margen de contribución anual por cliente --- ingreso menos los costos variables
    directamente atribuibles a atender a ese cliente. Usar margen, no ingreso; la fórmula
    mide el valor económico que acumula la empresa, no la facturación.
  \item[$r$] La tasa de descuento de la empresa (WACC o costo de capital propio) --- el
    costo de oportunidad del capital comprometido en atender clientes. Para la SaaS de
    referencia, $r = \qty{11}{\percent}$.
  \item[$g$] La tasa de crecimiento anual esperada de ese margen, que captura tanto la
    expansión de precio como de volumen dentro de la cuenta. Debe ser estrictamente menor
    que $r$ o la fórmula devuelve un resultado sin sentido.
\end{description}

\textbf{Qué mide.} El valor presente de todo el margen de contribución futuro que se espera
que genere un solo cliente, descontado a la tasa de costo de capital de la empresa.

\textbf{Por qué importa.} El LTV establece el \emph{techo} de cuánto debería estar dispuesta
a pagar la empresa por adquirir un cliente. Es matemáticamente idéntico a un valor terminal
de flujo de caja descontado (\cref{eq:gordon}) con el cliente reemplazando a la empresa; la
derivación completa está en \cref{ch:customer-economics}.

\textbf{Referencia.} $\mathrm{LTV} \ge 3 \times \mathrm{CAC}$. Esto no es arbitrario: una
razón de 3$\times$ deja margen para el costo de atender a los clientes durante su vida
(soporte, gasto en retención, gestión de cuentas) sin erosionar la economía de adquisición.

\textbf{Trampa.} Sustituir ingreso en lugar de margen en $m$ sobreestima el LTV de forma
dramática --- un negocio con \qty{42}{\percent} de margen de contribución que comete este
error infla el LTV en un \qty{138}{\percent}. Igualmente, usar una tasa de crecimiento $g$
elevada que se acerque a $r$ puede hacer el LTV arbitrariamente grande e inutilizar la
métrica para presupuestar.

\begin{example}[LTV para la SaaS de referencia]
El ARPU es \money{600}/año; el margen de contribución es \qty{42}{\percent}, por lo que $m =
\money{252}$/cliente/año. Tasa de descuento $r = \qty{11}{\percent}$; crecimiento del margen
a largo plazo $g = \qty{3}{\percent}$.
\[
  \mathrm{LTV} = \frac{\money{252}}{0.11 - 0.03} = \frac{\money{252}}{0.08} = \money{3150}.
\]
Si se hubiera usado el ingreso en lugar del margen: $\money{600}/0.08 = \money{7500}$ --- una
sobreestimación del \qty{138}{\percent} que distorsionaría el presupuesto de adquisición por
el mismo factor.
\end{example}


\subsection{LTV:CAC --- Razón entre el valor del ciclo de vida y el costo de adquisición}

\textbf{Fórmula.}
\[
  \mathrm{LTV:CAC} = \frac{\mathrm{LTV}}{\mathrm{CAC}}
\]

\textbf{Qué significa cada término.}
\begin{description}
  \item[LTV] Valor del ciclo de vida del cliente tal como se define arriba --- valor presente
    del margen futuro, no del ingreso.
  \item[CAC] Costo de adquisición de clientes completamente cargado tal como se define
    arriba.
\end{description}

\textbf{Qué mide.} La eficiencia de capital de la adquisición: ¿cuántos dólares de valor de
cliente crea la empresa por cada dólar gastado en ganar ese cliente?

\textbf{Por qué importa.} LTV:CAC es la prueba de economía unitaria más citada en la
suscripción de capital de riesgo y de crecimiento. Una razón inferior a 1 significa que el
negocio pierde dinero con cada cliente antes de contar siquiera los costos fijos. La fórmula
está completamente derivada como \cref{eq:ltv-cac} en \cref{ch:customer-economics}.

\textbf{Referencia.}
\begin{itemize}
  \item $\ge 3$: saludable; la referencia institucional estándar.
  \item 3--5: rango normal para SaaS bien gestionadas.
  \item $> 5$: con frecuencia indica subinversión en crecimiento --- la empresa está dejando
    valor sobre la mesa al no adquirir clientes que podría atender de forma rentable.
  \item $< 1$: cada nuevo cliente destruye valor; corregir la economía unitaria antes de
    escalar.
\end{itemize}

\textbf{Trampa.} Un LTV:CAC superior a 5 no es simplemente «muy bueno» --- típicamente
significa que la empresa gasta demasiado poco en adquisición en relación con la oportunidad.
Tratar una razón alta como recompensa en lugar de señal para aumentar el gasto es un error
frecuente en la fase de crecimiento.

\begin{example}[LTV:CAC para la SaaS de referencia]
De los cálculos anteriores, LTV $= \money{3150}$ y CAC $= \money{600}$.
\[
  \mathrm{LTV:CAC} = \frac{\money{3150}}{\money{600}} = 5.25.
\]
Esta cifra supera el rango saludable de 3--5, lo que sugiere que la empresa debería probar
aumentar el gasto en adquisición --- siempre que el cliente incremental tenga el mismo perfil
de LTV que el promedio.
\end{example}


\subsection{CAC Payback --- Período de recuperación en meses}

\textbf{Fórmula.}
\[
  \text{CAC Payback} = \frac{\mathrm{CAC}}{\text{Monthly margin per customer}}
\]

\textbf{Qué significa cada término.}
\begin{description}
  \item[CAC] Costo de adquisición completamente cargado por cliente (\money{600} en la SaaS
    de referencia).
  \item[Monthly margin per customer] Margen de contribución mensual --- ARPU mensual
    multiplicado por la tasa de margen de contribución. Es el flujo de caja que genera el
    cliente cada mes para recuperar la inversión de adquisición.
\end{description}

\textbf{Qué mide.} Cuántos meses debe permanecer un cliente antes de que el margen
acumulado de ese cliente reembolse el costo de adquirirlo. Los períodos de recuperación más
cortos reducen el riesgo del ciclo de caja y los requerimientos de financiamiento de la
empresa.

\textbf{Por qué importa.} El período de recuperación traduce el LTV:CAC en una verificación
de la realidad del flujo de caja. Un LTV:CAC elevado es tranquilizador pero carece de
sentido si los clientes sufren \emph{churn} antes del mes de recuperación. También es el
principal determinante de cuánto capital necesita la empresa para sostener una tasa de
crecimiento dada; un negocio creciendo al \qty{20}{\percent} mensual con 24 meses de
período de recuperación es mucho más intensivo en capital que uno con 8 meses. La derivación
completa es \cref{eq:cac-payback} en \cref{ch:customer-economics}.

\textbf{Referencia.}
\begin{itemize}
  \item $< 12$ meses: saludable para SaaS enfocadas en pequeñas y medianas empresas.
  \item $< 18$ meses: aceptable para mercado medio.
  \item $< 24$ meses: en el límite para empresa; por encima de 24 meses se requieren datos
    sólidos de retención para justificar el ciclo de caja.
\end{itemize}

\textbf{Trampa.} La fórmula cuenta a todos los clientes en el denominador, incluidos los que
sufren \emph{churn} antes del mes de recuperación. Si el \emph{churn} mensual es elevado, el
cliente promedio nunca alcanza la recuperación, convirtiendo la cifra declarada en una ficción
optimista. Verificar la recuperación contra el inverso del \emph{churn} mensual (la vida
media del cliente en meses).

\begin{example}[CAC Payback para la SaaS de referencia]
ARPU mensual $= \money{50}$; margen de contribución $= \qty{42}{\percent}$; margen mensual
por cliente $= \money{50} \times 0.42 = \money{21}$. CAC $= \money{600}$.
\[
  \text{CAC Payback} = \frac{\money{600}}{\money{21}} \approx 29 \text{ meses}.
\]
Con un \emph{churn} mensual de \qty{0.65}{\percent}, la vida media del cliente es
aproximadamente $1/0.0065 \approx 154$ meses, por lo que casi todos los clientes sobreviven
hasta la recuperación. Sin embargo, la cifra de 29 meses supera el umbral empresarial de
24 meses --- señal para reducir el CAC o encontrar palancas de expansión del margen
(venta adicional, aumento de precio) que compriman el período de recuperación.
\end{example}


\subsection{Contribution Margin --- Margen de contribución}

\textbf{Fórmula.}
\[
  \text{Contribution Margin} = \text{Price} - \text{Variable cost per unit}
\]
Expresado como tasa: $\text{CM\%} = (\text{Price} - \text{Variable cost})/\text{Price}$.

\textbf{Qué significa cada término.}
\begin{description}
  \item[Price] Ingreso por unidad (por asiento de suscripción, por transacción, por artículo
    vendido).
  \item[Variable cost per unit] Costos que escalan directamente con cada unidad adicional
    atendida: procesamiento de pagos, incrementos de alojamiento en la nube, APIs de terceros
    por asiento, material directo y mano de obra para bienes físicos. Los costos fijos
    (alquiler, salarios base, licencias anuales) se \emph{excluyen}.
\end{description}

\textbf{Qué mide.} El flujo de caja que genera cada unidad adicional una vez cubiertos los
costos variables --- el fondo de dinero disponible para cubrir los costos fijos y
eventualmente generar beneficio.

\textbf{Por qué importa.} El margen de contribución es el motor económico de la unidad. Un
margen de contribución positivo significa que cada venta incremental acerca la empresa al
umbral de rentabilidad; uno negativo significa que vender más acelera las pérdidas
independientemente del volumen. Se incorpora directamente al LTV (el término $m$) y a los
cálculos del período de recuperación del CAC. La relación con el margen bruto y el beneficio
operativo se desarrolla en \cref{ch:business-models}.

\textbf{Referencia.} Positivo y creciente con la escala a medida que mejora el apalancamiento
de los costos fijos. Para software, \qty{60}{\percent}+ es típico a nivel de unidad antes de
los gastos generales de plataforma; para bienes físicos, márgenes por debajo del
\qty{30}{\percent} hacen que la aritmética de economía unitaria sea muy ajustada.

\textbf{Trampa.} Confundir el margen de contribución con el margen bruto. El margen bruto
deduce el costo de los bienes vendidos tal como se reporta en el estado de resultados (que
puede incluir partidas semifijas como la depreciación del equipo de producción); el margen
de contribución es un concepto de contabilidad de gestión que aísla únicamente los costos
verdaderamente variables. Convergen para negocios de software puro pero divergen
materialmente para cualquier producto con componentes físicos o gran infraestructura de
producción fija.

\begin{example}[Margen de contribución en la SaaS de referencia]
Precio de suscripción mensual: \money{50}. Costos variables por asiento: procesamiento de
pagos \money{1.50} + alojamiento incremental en la nube \money{1.75} + costos de API por
asiento \money{0.75} = \money{4.00} total.
\[
  \text{CM} = \money{50} - \money{4} = \money{46}; \quad
  \text{CM\%} = \frac{\money{46}}{\money{50}} = \qty{92}{\percent}.
\]
Espera --- ¿por qué el \emph{fact sheet} de referencia usa \qty{42}{\percent}? Porque la
SaaS de referencia atribuye una parte de los salarios de éxito del cliente y las herramientas
de incorporación como costos variables en sus cuentas de gestión (un tratamiento completamente
cargado legítimo), elevando el costo variable por asiento a \money{29}. La cifra del
\qty{42}{\percent} (\money{21}/mes) es el número que fluye hacia el LTV y el período de
recuperación del CAC a lo largo de este libro.
\end{example}


\subsection{Gross Margin --- Margen bruto}

\textbf{Fórmula.}
\[
  \text{Gross Margin} = \frac{\text{Revenue} - \text{COGS}}{\text{Revenue}}
\]

\textbf{Qué significa cada término.}
\begin{description}
  \item[Revenue] Ingreso total reconocido en el período.
  \item[COGS (Cost of Goods Sold)] Los costos directos de entregar el producto o servicio
    tal como se reportan bajo GAAP/IFRS: mano de obra directa, materiales directos,
    alojamiento, licencias de terceros integradas en el producto, amortización de costos de
    desarrollo capitalizados directamente vinculados a la entrega de ingresos.
  \item[Gross Margin (como razón)] La proporción de cada dólar de ingreso que queda tras los
    costos directos de entrega; el punto de partida para todo análisis de rentabilidad
    posterior.
\end{description}

\textbf{Qué mide.} La rentabilidad estructural de la oferta principal de la empresa antes de
los gastos operativos --- ventas, \emph{marketing}, I+D y G\&A. Responde: ¿genera el
producto suficiente excedente para sostener la estructura operativa del negocio?

\textbf{Por qué importa.} El margen bruto es el techo estructural de la rentabilidad
operativa. Una empresa con \qty{40}{\percent} de margen bruto no puede, aritméticamente,
alcanzar un margen EBITDA del \qty{45}{\percent} --- ninguna eficiencia operativa supera ese
techo. Para los inversores en SaaS, el margen bruto señala escalabilidad: márgenes altos
($>\qty{70}{\percent}$) indican que el negocio puede crecer en ingresos sin aumentar
proporcionalmente los costos de entrega.

\textbf{Referencia.}
\begin{itemize}
  \item $>\qty{70}{\percent}$: rango saludable típico para software puro.
  \item $>\qty{80}{\percent}$: élite; característico de SaaS empresarial de alto valor.
  \item \qty{40}--\qty{60}{\percent}: común para software con componentes significativos de
    servicios gestionados o implementación.
  \item $<\qty{30}{\percent}$: típico para negocios de bienes físicos o \emph{marketplaces}
    donde el apalancamiento a nivel de unidad es menor.
\end{itemize}

\textbf{Trampa.} Enterrar los costos de alojamiento, los salarios del equipo de soporte o
los honorarios de APIs de terceros en los gastos operativos (en lugar de en COGS) embellece
el margen bruto en el estado de resultados pero engaña a cualquiera que modele la
escalabilidad. Esta reclasificación es una preparación estándar previa a la captación de
fondos en algunas \emph{startups}; un analista escéptico debería reconstruir el COGS a
partir del detalle de costos, no aceptar la cifra titular.

\begin{example}[Margen bruto en la SaaS de referencia]
ARR $= \money{4000000}$ (MRR \money{340000} $\times$ 12 aproximadamente). COGS reportado:
alojamiento \money{480000} + APIs de terceros \money{240000} + amortización de desarrollo
capitalizado \money{120000} = \money{840000}.
\[
  \text{Gross Margin} = \frac{\money{4000000} - \money{840000}}{\money{4000000}}
    = \frac{\money{3160000}}{\money{4000000}} = \qty{79}{\percent}.
\]
Si los salarios del equipo de soporte (\money{400000}) también se clasificaran correctamente
en COGS, el margen bruto caería a $(\money{4000000} - \money{1240000})/\money{4000000} =
\qty{69}{\percent}$ --- aún saludable pero por debajo del umbral de élite, cambiando el
relato para el inversor.
\end{example}


%% ============================================================
\section{SaaS e ingresos recurrentes}

Las métricas de ingresos recurrentes miden la salud del motor de suscripción. Son
indicadores adelantados: reflejan lo que hace la base de clientes hoy y predicen lo que
mostrará el estado de resultados en doce meses. Un fundador que observa únicamente los
ingresos sin estas métricas está mirando por el espejo retrovisor.

\subsection{MRR --- Monthly Recurring Revenue}

\textbf{Fórmula.}
\[
  \mathrm{MRR} = \sum_{\text{active subscriptions}} \text{monthly subscription value}
\]

\textbf{Qué significa cada término.}
\begin{description}
  \item[Monthly subscription value] El cargo recurrente normalizado por suscripción por mes.
    Los contratos anuales se dividen entre 12; los contratos plurianuales, entre el total de
    meses. Los honorarios por uso, los ingresos puntuales por servicios profesionales y las
    ventas de hardware se excluyen.
  \item[$\sum$ sobre suscripciones activas] Toda suscripción que esté activa (superada la
    prueba, sin \emph{churn}, sin pausa) en la fecha de medición.
\end{description}

\textbf{Qué mide.} La tasa de ejecución mensual predecible del negocio a partir de los
contratos de suscripción. Es el pulso de un negocio de ingresos recurrentes, actualizado
mensualmente a medida que se añaden nuevos clientes, los clientes existentes se expanden o
contraen y los clientes que sufren \emph{churn} salen.

\textbf{Por qué importa.} El MRR es la base de todo modelo financiero SaaS: ARR, NRR, burn
multiple y la Regla de 40 se derivan todos de él. Los inversores utilizan la tasa de
crecimiento del MRR --- no el crecimiento de ingresos --- como la señal de crecimiento
primaria para las empresas SaaS previas a la rentabilidad. El modelo de cascada completo del
MRR (nuevo + expansión $-$ contracción $-$ \emph{churn}) se deriva como \cref{eq:mrr-growth}
en \cref{ch:business-models}.

\textbf{Referencia.} La tasa de crecimiento importa más que el nivel. Un MRR de \money{340000}
creciendo al \qty{15}{\percent}/mes es fundamentalmente diferente de un MRR de \money{5000000}
creciendo al \qty{1}{\percent}/mes, aunque el segundo sea mayor en volumen.

\textbf{Trampa.} Contar ingresos no recurrentes --- honorarios de implementación puntuales,
excedentes de uso variable o compromisos de servicios profesionales --- como MRR. Esto infla
la métrica y, peor aún, crea una suavidad falsa: cuando esos proyectos puntuales terminan,
el MRR reportado cae y el negocio parece estar sufriendo \emph{churn}.

\begin{example}[Descomposición del MRR para la SaaS de referencia]
La empresa tiene \num{6700} clientes pagantes a \money{50}/mes de promedio.
\[
  \mathrm{MRR} = \num{6700} \times \money{50} = \money{335000}.
\]
El \emph{fact sheet} de referencia usa \money{340000} MRR, reflejando el recuento real de
clientes de \num{6800} en la fecha de medición (crecimiento en el cohorte respecto a los
\num{6700} del mes anterior). Los honorarios de implementación puntuales de \money{45000}
de este mes se excluyen; incluirlos sobreestimaría el MRR en un \qty{13}{\percent}.
\end{example}


\subsection{ARR --- Annual Recurring Revenue}

\textbf{Fórmula.}
\[
  \mathrm{ARR} = \mathrm{MRR} \times 12
\]

\textbf{Qué significa cada término.}
\begin{description}
  \item[MRR] Ingreso recurrente mensual tal como se define arriba --- solo suscripciones.
  \item[$\times 12$] La anualización supone que la tasa de ejecución actual se mantiene
    durante doce meses. Es una \emph{instantánea}, no una previsión.
\end{description}

\textbf{Qué mide.} La tasa de ejecución recurrente anualizada del negocio en un momento
dado. El ARR es la unidad estándar para los múltiplos de valoración SaaS (múltiplo de ARR,
derivado como \cref{eq:arr-multiple} en \cref{ch:business-models}) y para la comparación del
crecimiento año contra año.

\textbf{Por qué importa.} El ARR convierte la instantánea del MRR al lenguaje que usan los
inversores y adquirentes para los comparables. En el nivel de \money{4000000} ARR, la SaaS
de referencia está en el rango en el que los inversores institucionales de crecimiento
típicamente inician una diligencia debida seria; por debajo de \money{1000000} ARR, la
mayoría de los fondos institucionales de Serie A consideran que el conjunto de datos es
demasiado pequeño para extrapolar.

\textbf{Trampa.} Anualizar ingresos no recurrentes. Si un negocio factura \money{500000} en
consultoría puntual este mes y reporta un ARR de \money{6000000} multiplicando el ingreso
total por 12, ese «ARR» desaparecerá el mes siguiente. Los inversores que descubren el error
revalúan el negocio a su base recurrente real.

\begin{example}[ARR para la SaaS de referencia]
MRR $= \money{340000}$.
\[
  \mathrm{ARR} = \money{340000} \times 12 = \money{4080000}.
\]
El \emph{fact sheet} de referencia usa \money{4000000} ARR, representando una línea base de
MRR de \money{333333} --- el redondeo menor refleja la cifra normalizada usada para
modelización. Si la empresa hubiera incluido los honorarios de implementación del mes
anterior, el ARR reportado habría sido \money{4620000} --- una sobreestimación del
\qty{15}{\percent} sin sustancia recurrente.
\end{example}


\subsection{Logo Churn --- Tasa de cancelación de clientes}

\textbf{Fórmula.}
\[
  \text{Logo Churn} = \frac{\text{Customers lost in period}}{\text{Customers at start of period}}
\]

\textbf{Qué significa cada término.}
\begin{description}
  \item[Customers lost in period] Suscripciones que cancelaron o no renovaron durante el
    período de medición (típicamente mensual o anual). Excluye las reducciones de plan
    (que son contracción, no \emph{churn}).
  \item[Customers at start of period] La base de la cohorte. Medir al inicio del período,
    antes de las nuevas incorporaciones.
\end{description}

\textbf{Qué mide.} La tasa a la que el negocio pierde clientes por número, independientemente
del tamaño de sus ingresos.

\textbf{Por qué importa.} El Logo Churn es la señal de desgaste más intuitiva: si el
\qty{5}{\percent} de los clientes se va cada mes, toda la base rota en aproximadamente
20 meses. Con un \emph{churn} mensual del \qty{0.65}{\percent}, la vida media del cliente
de la SaaS de referencia es $1/0.0065 \approx 154$ meses --- más de 12 años --- un relato
de retención muy sólido. El \emph{churn} se incorpora directamente al LTV: cada mejora de
un punto porcentual en el \emph{churn} mensual extiende la vida del cliente y eleva el LTV
de forma no lineal a través del denominador de la perpetuidad. La derivación de retención es
\cref{eq:retention} en \cref{ch:customer-economics}.

\textbf{Referencia.}
\begin{itemize}
  \item $<\qty{1}{\percent}$/mes (\qty{12}{\percent}/año): aceptable para SaaS enfocadas en
    pequeñas y medianas empresas.
  \item $<\qty{0.5}{\percent}$/mes (\qty{6}{\percent}/año): bueno para mercado medio.
  \item $<\qty{0.2}{\percent}$/mes: retención empresarial de élite.
  \item Los contratos anuales suprimen naturalmente las cifras de \emph{churn} mensual;
    comparar sobre una base consistente.
\end{itemize}

\textbf{Trampa.} Ignorar que los clientes perdidos pueden ser cuentas pequeñas. Si los diez
clientes más grandes (cada uno a \money{10000}/mes) se van pero se añaden doscientas cuentas
pequeñas (cada una a \money{50}/mes), el Logo Churn parece saludable mientras el \emph{churn}
por ingresos es catastrófico. Siempre combinar el Logo Churn con la retención neta de ingresos
(\cref{sec:saas-mrr} a continuación) para obtener el panorama completo.

\begin{example}[Logo Churn en la SaaS de referencia]
Recuento de clientes al inicio: \num{6700}. Cancelaciones en el mes: 44.
\[
  \text{Logo Churn} = \frac{44}{\num{6700}} = \qty{0.657}{\percent}/\text{mes}.
\]
El \emph{fact sheet} de referencia usa \qty{0.65}{\percent}/mes. A esta tasa, la vida media
del cliente $\approx 154$ meses; el período de recuperación del CAC de 29 meses queda bien
dentro de esa ventana.
\end{example}


\subsection{NRR --- Net Revenue Retention}

\textbf{Fórmula.}
\[
  \mathrm{NRR} = \frac{\text{MRR}_\text{start} + \text{Expansion MRR} - \text{Contraction MRR} - \text{Churned MRR}}{\text{MRR}_\text{start}}
\]

\textbf{Qué significa cada término.}
\begin{description}
  \item[MRR$_\text{start}$] Ingreso recurrente mensual de la base de clientes existentes al
    inicio del período --- sin incluir nuevos clientes.
  \item[Expansion MRR] MRR adicional de clientes existentes mediante ventas adicionales,
    ventas cruzadas o expansión de asientos.
  \item[Contraction MRR] MRR perdido de clientes existentes que redujeron su plan sin
    cancelar.
  \item[Churned MRR] MRR de clientes que cancelaron completamente.
\end{description}

\textbf{Qué mide.} El porcentaje de los ingresos iniciales que el negocio retiene de su
base de clientes existentes tras doce meses, neto de expansiones, contracciones y
cancelaciones.

\textbf{Por qué importa.} Un NRR superior al \qty{100}{\percent} significa que la base de
clientes existente hace crecer los ingresos incluso con cero adquisición de nuevos clientes.
Esta es la característica definitoria de un negocio SaaS de capitalización compuesta: el
crecimiento no requiere llenar perpetuamente un cubo con fugas. Con
$\mathrm{NRR} = \qty{120}{\percent}$, una empresa que comience con \money{4000000} ARR
alcanzaría \money{4800000} ARR solo desde su base actual en doce meses. Véase
\cref{ch:business-models} para la mecánica completa.

\textbf{Referencia.}
\begin{itemize}
  \item $>\qty{100}{\percent}$: expansión neta --- la base existente se hace crecer a sí
    misma.
  \item $>\qty{110}{\percent}$: bueno; referencia de SaaS de mercado medio.
  \item $>\qty{120}{\percent}$: élite; característico de las mejores empresas SaaS
    empresariales.
  \item $<\qty{100}{\percent}$: la base se está contrayendo incluso antes del Logo Churn;
    una señal de alerta grave que requiere atención inmediata.
\end{itemize}

\textbf{Trampa.} Un NRR elevado puede enmascarar un Logo Churn intenso. Si el negocio está
perdiendo muchos clientes pequeños pero los pocos clientes grandes se están expandiendo, el
NRR se ve sólido mientras la base de clientes se concentra en un número reducido de cuentas
grandes --- una fragilidad que es invisible en la cifra de NRR por sí sola.

\begin{example}[NRR para la SaaS de referencia]
MRR inicial de los \num{6700} clientes existentes: \money{335000}. Expansión (venta adicional
a nivel superior): \money{12000}. Contracción (reducciones de plan): \money{3500}. MRR por
\emph{churn} (44 clientes $\times$ \money{50}): \money{2200}.
\[
  \mathrm{NRR} = \frac{\money{335000} + \money{12000} - \money{3500} - \money{2200}}
    {\money{335000}}
  = \frac{\money{341300}}{\money{335000}} \approx \qty{102}{\percent}.
\]
El negocio está en expansión neta --- aunque modestamente. Duplicar el movimiento de
expansión (llevando el NRR al \qty{104}{\percent}) añadiría aproximadamente \money{13000}/mes
en crecimiento orgánico de ARR solo desde la base existente.
\end{example}


\subsection{Rule of 40 --- Regla de 40}

\textbf{Fórmula.}
\[
  \text{Rule of 40} = \text{Revenue growth rate (\%)} + \text{Profit margin (\%)}
\]

\textbf{Qué significa cada término.}
\begin{description}
  \item[Revenue growth rate] Crecimiento de ARR o ingresos año contra año expresado como
    porcentaje.
  \item[Profit margin] Típicamente margen EBITDA o, para empresas en etapas más tempranas,
    margen de flujo de caja libre. Los dos no son intercambiables; aclarar cuál se usa antes
    de comparar entre empresas.
\end{description}

\textbf{Qué mide.} El equilibrio entre crecimiento y rentabilidad. Una empresa que gasta de
forma agresiva para crecer (margen de beneficio negativo) puede seguir satisfaciendo la Regla
de 40 si su tasa de crecimiento es suficientemente alta, y viceversa.

\textbf{Por qué importa.} La Regla de 40 es la heurística principal que los inversores en
SaaS usan para evaluar la compensación entre crecimiento y eficiencia a escala. Reconoce que
la combinación «correcta» de crecimiento y margen no es fija --- una empresa con
\qty{60}{\percent} de crecimiento a $-\qty{20}{\percent}$ de margen y una con
\qty{10}{\percent} de crecimiento a \qty{30}{\percent} de margen obtienen la misma
puntuación, y ambas son aceptables en diferentes fases. La derivación formal es
\cref{eq:rule40} en \cref{ch:business-models}.

\textbf{Referencia.} $\ge 40$: el umbral en el que los inversores típicamente asignan
múltiplos premium. Por encima de 60 es de élite. Por debajo de 40 plantea preguntas sobre
la sostenibilidad del perfil crecimiento-margen.

\textbf{Trampa.} Manipular la puntuación con partidas de margen puntuales (una gran venta de
activos, un crédito por reestructuración, una recuperación de ingresos diferidos) que inflan
el margen de beneficio en un solo período sin cambiar el negocio subyacente. Examinar siempre
los últimos doce meses en múltiples trimestres para suavizar los elementos excepcionales.

\begin{example}[Regla de 40 para la SaaS de referencia]
El ARR creció de \money{3000000} a \money{4000000} en doce meses: tasa de crecimiento
$= \qty{33}{\percent}$. Margen EBITDA del año: \qty{9}{\percent} (próximo al NOPAT,
consistente con \money{600000} NOPAT sobre \money{4000000} de ingresos, margen NOPAT del
\qty{15}{\percent} aproximado al \qty{9}{\percent} EBITDA tras ajustes).
\[
  \text{Rule of 40} = 33 + 9 = 42.
\]
El negocio supera el umbral --- por poco. Para alcanzar 50 sin acelerar el crecimiento, la
empresa necesitaría expandir el margen EBITDA al \qty{17}{\percent}, lo que al ingreso actual
significa \money{680000} adicionales de beneficio operativo --- alcanzable únicamente mediante
poder de fijación de precios o apalancamiento de costos fijos, no mediante volumen incremental
a los márgenes actuales.
\end{example}


\subsection{Burn Multiple --- Múltiplo de quema}

\textbf{Fórmula.}
\[
  \text{Burn Multiple} = \frac{\text{Net cash burn}}{\text{Net new ARR}}
\]

\textbf{Qué significa cada término.}
\begin{description}
  \item[Net cash burn] Efectivo consumido por el negocio en el período --- salida de caja
    operativa neta de cualquier entrada de caja procedente de operaciones (cobros, no
    captación de fondos).
  \item[Net new ARR] El ARR incremental añadido en el mismo período: ARR de nuevos clientes
    + ARR de expansión $-$ ARR por \emph{churn} $-$ ARR por contracción. Es el ARR que la
    quema «compró».
\end{description}

\textbf{Qué mide.} Cuántos dólares de caja quema el negocio para generar cada dólar de ARR
neto nuevo. Es una métrica de eficiencia de capital para el crecimiento --- cuanto menor,
mejor.

\textbf{Por qué importa.} El Burn Multiple conecta el estado de resultados con el balance:
un negocio puede estar haciendo crecer el ARR rápidamente mientras quema caja a una tasa que
agotará la \emph{runway} antes del siguiente evento de financiación. Con \$1 de quema por
\$1 de ARR nuevo, la empresa paga un dólar de capital por un dólar de ingreso futuro ---
eficiente. Con \$5 por \$1, gasta \$5 para adquirir \$1 de ARR, lo que implica o bien un
movimiento de crecimiento de muy baja calidad o un problema estructural de economía unitaria.
La definición formal es \cref{eq:burn-multiple} en \cref{ch:business-models}.

\textbf{Referencia.}
\begin{itemize}
  \item $< 1$: excelente eficiencia de capital; el negocio está cerca de ser positivo en
    flujo de caja desde la perspectiva del crecimiento.
  \item $< 1.5$: bueno; la mayoría de SaaS en etapa de crecimiento bien gestionadas.
  \item $1.5$--$3$: aceptable en una fase de expansión territorial deliberada con NRR sólido.
  \item $> 3$: preocupante; sugiere que el crecimiento es costoso en relación con su calidad.
\end{itemize}

\textbf{Trampa.} La métrica empeora en fases de expansión territorial deliberada donde la
empresa está invirtiendo por delante de los ingresos en cobertura de mercado, ingeniería o
plantilla de éxito del cliente. El contexto importa: un múltiplo de quema de \$2.5 junto con
un NRR del \qty{150}{\percent} es una situación fundamentalmente diferente a un múltiplo de
\$2.5 con un NRR del \qty{90}{\percent}.

\begin{example}[Burn Multiple para la SaaS de referencia]
Quema de caja neta en el T1: \money{450000}. ARR neto nuevo en el T1: nuevos clientes
(\num{200} clientes $\times$ \money{600}/año) = \money{120000}; expansión \money{48000};
pérdida por \emph{churn} \money{26400}; ARR neto nuevo $= \money{141600}$.
\[
  \text{Burn Multiple} = \frac{\money{450000}}{\money{141600}} \approx 3.18.
\]
Esta cifra supera el umbral «bueno». Descomponiendo la quema se revelan \money{120000} en
contrataciones de I+D por delante del lanzamiento de una funcionalidad empresarial --- si ese
lanzamiento eleva el NRR en 4 puntos porcentuales, el múltiplo efectivo a futuro caerá a
aproximadamente 1.8.
\end{example}


%% ============================================================
\section{Embudo y marketing}

Las métricas de \emph{funnel} rastrean la conversión de atención en ingresos en cada etapa
del recorrido del cliente. Ninguna métrica de \emph{funnel} debe optimizarse de forma
aislada --- una tasa de conversión que sube mientras cae el valor medio de pedido puede
representar una pérdida neta; un CPA que cae pero trae prospectos de menor calidad erosiona
silenciosamente el LTV.

\subsection{CTR --- Click-Through Rate}

\textbf{Fórmula.}
\[
  \mathrm{CTR} = \frac{\text{Clicks}}{\text{Impressions}}
\]

\textbf{Qué significa cada término.}
\begin{description}
  \item[Clicks] El número de veces que un usuario hizo clic en el anuncio o enlace durante
    el período de medición.
  \item[Impressions] El número de veces que el anuncio o enlace se mostró a los usuarios
    (incluidos aquellos que no hicieron clic).
\end{description}

\textbf{Qué mide.} La proporción de personas que vieron el anuncio o contenido y optaron por
interactuar. El CTR es una señal directa de la resonancia creativa y de segmentación ---
si el mensaje es suficientemente relevante para interrumpir la atención.

\textbf{Por qué importa.} El CTR es un diagnóstico de la parte alta del \emph{funnel}. Un
CTR alto indica que el creativo, el titular o la oferta es convincente \emph{para la
audiencia que se le está mostrando}. No indica la calidad de la conversión posterior. Monitorear
el CTR junto a la tasa de conversión (más adelante) revela si el problema de una campaña
está en la etapa de atención o en la etapa de aterrizaje/oferta. El CTR en plataformas de
pago también alimenta el Índice de Calidad y la mecánica de eCPC, desarrollados en
\cref{ch:paid-acquisition-platforms}.

\textbf{Referencia.} Varía ampliamente según el canal y el formato; los anuncios de búsqueda
suelen rondar el \qty{2}{\percent}--\qty{6}{\percent}, los anuncios de display el
\qty{0.1}{\percent}--\qty{0.3}{\percent}, los enlaces de email el
\qty{1}{\percent}--\qty{3}{\percent}. Comparar únicamente dentro de los mismos canales y
contra la propia línea base histórica.

\textbf{Trampa.} Un CTR alto con una tasa de conversión posterior baja significa que el
anuncio está atrayendo a la audiencia equivocada --- o prometiendo algo que la página de
destino no entrega. Optimizar solo para CTR (escribiendo titulares engañosos o
sensacionalistas) es una manera segura de destruir la eficiencia del CAC mientras se alcanza
una métrica de vanidad.

\begin{example}[CTR en la SaaS de referencia]
Una campaña de LinkedIn dirigida a directores de operaciones: \num{240000} impresiones;
\num{4800} clics.
\[
  \mathrm{CTR} = \frac{4800}{240000} = \qty{2}{\percent}.
\]
De esos \num{4800} clics, \num{96} se convirtieron en registros de prueba (tasa de conversión
\qty{2}{\percent}) y \num{20} en clientes pagantes. El CTR parece saludable, pero la tasa
de clic a prueba revela un problema en la página de destino --- la oferta resuena en el
feed pero no convierte una vez que se llega a ella.
\end{example}


\subsection{CPC --- Cost Per Click}

\textbf{Fórmula.}
\[
  \mathrm{CPC} = \frac{\text{Total ad spend}}{\text{Total clicks}}
\]

\textbf{Qué significa cada término.}
\begin{description}
  \item[Total ad spend] El importe pagado a la plataforma por la campaña durante el período
    de medición.
  \item[Total clicks] El número de clics entregados por la plataforma en el mismo período.
\end{description}

\textbf{Qué mide.} El precio efectivo que paga el negocio por una sola visita desde un
anuncio de pago. El CPC lo establece la subasta de la plataforma y el Índice de Calidad del
anunciante (el mecanismo \cref{eq:ecpc} de \cref{ch:paid-acquisition-platforms}).

\textbf{Por qué importa.} El CPC es el costo de entrada para toda la economía del tráfico
de pago. Con una tasa de conversión y un valor medio de pedido dados, el CPC determina si
el tráfico de pago es rentable. La relación es: $\mathrm{CPA} = \mathrm{CPC} /
\mathrm{CVR}$, por lo que un CPC elevado requiere una tasa de conversión alta o un CLV alto
para justificarse.

\textbf{Referencia.} No existe una referencia universal; comparar contra el techo de CPA
implícito por el LTV. Google Search en SaaS B2B suele rondar \money{15}--\money{60} por clic
para términos competitivos; LinkedIn puede rondar \money{8}--\money{25} por clic.

\textbf{Trampa.} Optimizar para el CPC más bajo posible en lugar del CPA más bajo. El CPC
puede reducirse pujando por palabras clave más baratas y de menor intención --- pero si esos
clics convierten a una cuarta parte de la tasa, el CPA sube aunque el CPC baje. Evaluar
siempre el CPC en el contexto de la conversión posterior.

\begin{example}[CPC en la SaaS de referencia]
Campaña de LinkedIn: gasto \money{12000}; \num{4800} clics.
\[
  \mathrm{CPC} = \frac{\money{12000}}{4800} = \money{2.50}.
\]
Con una tasa de conversión de clic a pagante de $20/4800 = \qty{0.42}{\percent}$:
$\mathrm{CPA} = \money{2.50}/0.0042 \approx \money{595}$ --- casi exactamente el CAC de
\money{600}. Esta campaña opera en el umbral de rentabilidad de la economía unitaria.
\end{example}


\subsection{CPM --- Cost Per Thousand Impressions}

\textbf{Fórmula.}
\[
  \mathrm{CPM} = \frac{\text{Total ad spend}}{\text{Total impressions}} \times 1000
\]

\textbf{Qué significa cada término.}
\begin{description}
  \item[Total ad spend] Importe pagado a la plataforma durante el período de medición.
  \item[Total impressions] Número de veces que el anuncio se mostró (sirvió) a los usuarios.
  \item[$\times 1000$] Normaliza la cifra a una base por mil --- la unidad estándar de la
    industria para comparar los costos de alcance entre canales.
\end{description}

\textbf{Qué mide.} El costo de llegar a mil espectadores con una impresión de anuncio. El
CPM es la unidad principal de compra y fijación de precios para campañas de
marca/notoriedad donde el alcance y la frecuencia importan más que la conversión directa.

\textbf{Por qué importa.} El CPM gobierna el costo de construir notoriedad en la parte alta
del \emph{funnel}. Cuando una marca hace publicidad en televisión, display, audio en
streaming o display programático, el CPM determina cuánto alcance de audiencia compra un
presupuesto dado. En \emph{marketing} de respuesta directa, el CPM es un derivado del CTR
y el CPC más que un objetivo de optimización primario.

\textbf{Referencia.} Display programático: CPM de \money{2}--\money{8}. Feeds de redes
sociales: CPM de \money{6}--\money{20}. LinkedIn: CPM de \money{30}--\money{80} (la
segmentación B2B premium conlleva una gran prima). TV conectada: CPM de
\money{20}--\money{50}.

\textbf{Trampa.} Optimizar para un CPM barato en una audiencia de escaso valor. Una campaña
con CPM de \money{1} que llega a usuarios sin interés en el producto desperdicia cada
impresión; una campaña de LinkedIn con CPM de \money{40} que llega exactamente al perfil de
quien toma la decisión puede ser dramáticamente más eficiente en términos de costo por
resultado.

\begin{example}[CPM en la SaaS de referencia]
Una campaña de retargeting programático en display: gasto \money{3000}; impresiones
\num{500000}.
\[
  \mathrm{CPM} = \frac{\money{3000}}{500000} \times 1000 = \money{6}.
\]
Clics: \num{750} (CTR \qty{0.15}{\percent}); conversiones a pago: 3.
$\mathrm{CPA} = \money{3000}/3 = \money{1000}$ --- significativamente por encima del CAC.
El CPM barato está comprando impresiones de escaso valor; la audiencia de retargeting
necesita afinar su definición.
\end{example}


\subsection{CPA / CPL --- Cost Per Acquisition / Cost Per Lead}

\textbf{Fórmula.}
\[
  \mathrm{CPA} = \frac{\text{Ad spend}}{\text{Conversions}}
\]

\textbf{Qué significa cada término.}
\begin{description}
  \item[Ad spend] El costo de medios de la campaña (excluye los costos de producción a
    menos que se especifique como CPA completamente cargado).
  \item[Conversions] La acción específica que se mide --- un cliente pagante (CPA), un
    prospecto calificado por ventas (CPL), un registro de prueba, un formulario enviado,
    dependiendo del contexto de medición. La definición debe ser fija antes de comparar el
    CPA entre campañas.
\end{description}

\textbf{Qué mide.} El costo promedio de impulsar una acción objetivo desde los medios de
pago. El CPA es la métrica de eficiencia primaria del profesional de \emph{marketing} de
respuesta directa, situada un nivel por encima de CTR/CPC y directamente comparable al CAC
cuando la acción es un cliente pagante.

\textbf{Por qué importa.} El CPA es el nexo entre el gasto en medios y la economía unitaria:
si $\mathrm{CPA} < \mathrm{LTV}/3$, el canal es económicamente viable al LTV:CAC objetivo.
El techo para el CPA se deriva del LTV en \cref{ch:paid-acquisition-platforms}.

\textbf{Referencia.} El techo implícito por el LTV: $\mathrm{CPA}_{\max} = \mathrm{LTV} / 3$
para un LTV:CAC de 3. Para la SaaS de referencia,
$\mathrm{CPA}_{\max} = \money{3150}/3 = \money{1050}$ por cliente pagante.

\textbf{Trampa.} Definir las conversiones de manera demasiado laxa para embellecer la
métrica. Contar formularios enviados (CPL a \money{20}) en lugar de prospectos calificados
por ventas (CPSQL a \money{200}) puede hacer que una campaña parezca 10$\times$ más eficiente
de lo que realmente es a nivel de ingresos. Cada cifra de CPA requiere una definición clara
y estable del evento de conversión.

\begin{example}[CPA para la SaaS de referencia]
Campaña de Google Search para el mes: gasto \money{52000}; nuevos clientes pagantes: 87.
\[
  \mathrm{CPA} = \frac{\money{52000}}{87} \approx \money{598}.
\]
Esta cifra queda justo por debajo del CAC combinado de \money{600}, confirmando que el canal
de Google Search contribuye aproximadamente a la eficiencia media de adquisición. El techo
implícito por el LTV de \money{1050} significa que hay margen significativo para aumentar las
pujas antes de que el canal se vuelva no económico.
\end{example}


\subsection{ROAS --- Return on Ad Spend}

\textbf{Fórmula.}
\[
  \mathrm{ROAS} = \frac{\text{Conversion value (revenue)}}{\text{Ad spend}}
\]

\textbf{Qué significa cada término.}
\begin{description}
  \item[Conversion value] Ingresos atribuidos a las conversiones impulsadas por la campaña,
    tal como reporta la plataforma publicitaria. El modelo de atribución de la plataforma
    (típicamente último clic) determina qué ingresos se acreditan.
  \item[Ad spend] El costo de medios de la campaña durante el período de medición.
\end{description}

\textbf{Qué mide.} Ingresos devueltos por dólar de gasto publicitario, medidos por la
atribución propia de la plataforma. El ROAS es la métrica estándar de eficiencia para
campañas de comercio electrónico y respuesta directa donde la plataforma puede rastrear los
ingresos directamente.

\textbf{Por qué importa.} El ROAS traduce el costo de medios en un múltiplo de ingresos,
permitiendo una comparación rápida entre campañas y canales. Se conecta directamente con la
rentabilidad: un ROAS de 3 con un margen bruto del \qty{40}{\percent} significa que la
campaña está en el umbral de rentabilidad del beneficio bruto ($3 \times 0.40 = 1.20$ de
beneficio bruto por dólar gastado, menos el dólar de gasto mismo deja \money{0.20} de
cobertura de margen bruto, que debe absorber los costos fijos).

\textbf{Referencia.} Para comercio electrónico: 3--5$\times$ es ampliamente aceptable;
5--8$\times$ es fuerte. El objetivo de ROAS correcto debe derivarse del margen bruto:
$\mathrm{ROAS}_{\min} = 1/\text{margen bruto}$ para el umbral de rentabilidad del beneficio
bruto. Todas las cifras de ROAS reportadas por las plataformas están sujetas a inflación por
atribución (\cref{ch:attribution-and-measurement}).

\textbf{Trampa.} El ROAS reportado por la plataforma acredita sistemáticamente en exceso a
la propia plataforma. El modelo de atribución de cada canal reclama crédito por conversiones
que también aparecen en los reportes de otros canales (doble conteo), y los créditos por
visualización inflan el ROAS en canales con muchas impresiones. Usar el ROAS como herramienta
de gestión de campañas, no como comparación de eficiencia entre canales; usar el iROAS para
la medición causal.

\begin{example}[ROAS para un caso ilustrativo de consumo directo]
Una \textit{marca de suplementos de consumo} (no SaaS, a título ilustrativo) lanza una
campaña Meta Advantage+: gasto \money{25000}; ingresos reportados por la plataforma:
\money{125000}; ROAS $= 5\times$. Una prueba de incrementalidad (retención geográfica)
revela los ingresos incrementales reales: \money{60000}; iROAS $= 2.4\times$ --- apenas por
encima del umbral de rentabilidad con el margen bruto del \qty{45}{\percent} de la empresa.
La plataforma sobrecreditó el ROAS en más de un \qty{100}{\percent}.
\end{example}


\subsection{MER --- Marketing Efficiency Ratio}

\textbf{Fórmula.}
\[
  \mathrm{MER} = \frac{\text{Total revenue}}{\text{Total marketing spend}}
\]

\textbf{Qué significa cada término.}
\begin{description}
  \item[Total revenue] Todos los ingresos del período --- no solo los ingresos atribuidos. La
    facturación total del negocio.
  \item[Total marketing spend] Todo el gasto de \emph{marketing} en todos los canales de
    pago en el mismo período.
\end{description}

\textbf{Qué mide.} La eficiencia combinada y global de la inversión en \emph{marketing}.
A diferencia del ROAS por canal (que depende de modelos de atribución), el MER es libre de
modelos: simplemente pregunta cuántos ingresos genera el negocio por dólar de gasto total
en \emph{marketing}.

\textbf{Por qué importa.} El MER esquiva completamente el debate de atribución. Cuando las
cifras de ROAS por canal de Meta, Google y TikTok suman más que los ingresos totales del
negocio (porque cada plataforma sobrecontabiliza), el MER proporciona la verificación
honesta. Un negocio con MER de \money{4} sabe que por cada dólar de gasto total en
\emph{marketing}, el negocio genera \money{4} de ingreso total --- independientemente de
qué canal lo «causó».

\textbf{Referencia.} Varía ampliamente según el modelo de negocio y la estructura de
márgenes. Para SaaS: un MER de 10--20$\times$ sobre el gasto de \emph{marketing} para nuevos
clientes es típico a escala (porque el LTV es alto y el CAC es bajo en relación con el ACV).
Para comercio electrónico con márgenes escasos, un MER de 3--6$\times$ puede ser el rango
viable.

\textbf{Trampa.} El MER oculta qué canal impulsó realmente los ingresos. Una empresa que
desactiva su canal de mejor desempeño y traslada el presupuesto a uno peor puede mantener el
MER a corto plazo (porque el capital de marca del buen canal sostiene los ingresos
brevemente) mientras degrada sistemáticamente el desempeño futuro. Usar el MER como
verificación de salud agregada, no como herramienta de asignación de canales.

\begin{example}[MER para la SaaS de referencia]
MRR total este mes: \money{340000} (todos los ingresos, nuevos y recurrentes). Gasto total
en \emph{marketing}: \money{120000} (el numerador del CAC combinado).
\[
  \mathrm{MER} = \frac{\money{340000}}{\money{120000}} \approx 2.83.
\]
Este MER es deliberadamente bajo porque divide el ingreso \emph{mensual} total entre el gasto
solo de adquisición (la mayor parte de los ingresos proviene de clientes existentes, no de
nuevas adquisiciones). En una base de solo ingresos por nuevos clientes (200 nuevos clientes
$\times$ \money{50}/mes $= \money{10000}$): $\mathrm{MER}_\text{adquisición} =
\money{10000}/\money{120000} = 0.08$ --- lo cual es esperable, dado que el modelo de negocio
es recurrente y el período de recuperación es de 29 meses.
\end{example}


\subsection{AOV --- Average Order Value}

\textbf{Fórmula.}
\[
  \mathrm{AOV} = \frac{\text{Total revenue}}{\text{Number of orders}}
\]

\textbf{Qué significa cada término.}
\begin{description}
  \item[Total revenue] Ingresos de todas las transacciones en el período de medición.
  \item[Number of orders] Número de transacciones de compra distintas en el mismo período.
\end{description}

\textbf{Qué mide.} El tamaño promedio de una sola transacción. El AOV es principalmente una
métrica de comercio minorista y electrónico; en SaaS el equivalente es el ARPU (ingreso
medio por usuario).

\textbf{Por qué importa.} El AOV determina directamente el CPA máximo viable: un canal que
atrae compradores con un AOV de \money{30} y un margen bruto del \qty{40}{\percent} tiene
como máximo \money{12} de beneficio bruto para financiar la adquisición. Elevar el AOV
mediante ventas adicionales, ventas cruzadas o agrupación de productos es una de las
intervenciones de mayor apalancamiento en el comercio electrónico porque mejora la
rentabilidad sin cambiar el volumen de tráfico ni la tasa de conversión.

\textbf{Referencia.} Completamente dependiente de la categoría. Las referencias de AOV solo
tienen sentido dentro de una misma vertical (p. ej., comercio electrónico de moda,
suscripciones de software B2B). La tendencia a lo largo del tiempo importa más que el nivel
absoluto.

\textbf{Trampa.} La media oculta una cola larga y sesgada. Un número reducido de pedidos
grandes puede elevar el AOV muy por encima del gasto típico del cliente, haciendo que el
promedio sea inútil para las decisiones de fijación de precios y agrupación. Examinar siempre
la mediana y la distribución, no solo la media.

\begin{example}[AOV para un caso ilustrativo de comercio electrónico]
Una \textit{marca de café de venta directa al consumidor} (ilustrativo): \num{4200} pedidos
en el mes; ingresos \money{189000}.
\[
  \mathrm{AOV} = \frac{\money{189000}}{4200} = \money{45}.
\]
La media es \money{45} pero el pedido mediano es \money{32} --- la brecha está impulsada por
200 pedidos de paquetes de suscripción con un promedio de \money{120}. Orientar la
publicidad hacia los compradores de paquetes eleva el AOV y mejora drásticamente la
viabilidad del CPA.
\end{example}


\subsection{Conversion Rate --- Tasa de conversión}

\textbf{Fórmula.}
\[
  \mathrm{CVR} = \frac{\text{Conversions}}{\text{Sessions (or visitors)}}
\]

\textbf{Qué significa cada término.}
\begin{description}
  \item[Conversions] La acción objetivo específica: una compra de pago, un registro de
    prueba, una demo reservada o cualquier otro evento definido. La acción debe coincidir con
    el propósito de la métrica.
  \item[Sessions] El número de sesiones de usuario distintas (visitas) a la página o etapa
    del \emph{funnel} medida en el mismo período.
\end{description}

\textbf{Qué mide.} La eficiencia de una etapa específica del \emph{funnel} --- la fracción
de visitantes que realizan la acción objetivo. El CVR puede medirse en cada etapa: sesión a
prueba, prueba a pago, demo a cierre, etc.

\textbf{Por qué importa.} El CVR vincula el volumen de tráfico con los resultados de
ingresos y, por lo tanto, vincula el CPC con el CPA ($\mathrm{CPA} = \mathrm{CPC}/\mathrm{CVR}$).
Una mejora del \qty{1}{\percent} en la tasa de conversión tiene el mismo efecto sobre el CPA
que una reducción del \qty{1}{\percent} en el CPC --- pero la mejora de la tasa de conversión
típicamente cuesta tiempo de I+D y diseño en lugar de gasto en medios. La cadena completa
del \emph{funnel} se formaliza como \cref{eq:funnel-chain} en \cref{ch:digital-funnel-and-crm}.

\textbf{Referencia.} Altamente dependiente del contexto. Sesión a prueba para SaaS B2B:
\qty{2}{\percent}--\qty{5}{\percent}. Prueba a pago: \qty{15}{\percent}--\qty{30}{\percent}
con incorporación activa. CVR de comercio electrónico a nivel de sitio:
\qty{1}{\percent}--\qty{4}{\percent}.

\textbf{Trampa.} El CVR a nivel de sitio no tiene casi ningún significado sin segmentación
por fuente de tráfico, dispositivo, página de destino e intención del usuario. Los visitantes
de búsqueda de marca convierten a 5--10$\times$ la tasa del tráfico de display en frío;
promediarlos produce un número que no describe con precisión a ninguno de los dos grupos y
distorsiona las decisiones de optimización en ambas direcciones.

\begin{example}[Tasa de conversión en la SaaS de referencia]
Del \emph{fact sheet} de referencia: \num{10000} sesiones mensuales $\to$ \num{800} pruebas
$\to$ \num{200} clientes pagantes.
\[
  \mathrm{CVR}_{\text{sesión a prueba}} = \frac{800}{10000} = \qty{8}{\percent}; \qquad
  \mathrm{CVR}_{\text{prueba a pago}} = \frac{200}{800} = \qty{25}{\percent}.
\]
La tasa de sesión a prueba es sólida (por encima de la referencia típica), lo que sugiere
una buena intención de SEO orgánico. La tasa de prueba a pago está en el extremo superior
del rango saludable, indicando una incorporación efectiva. La restricción es el volumen de
tráfico en la parte alta del \emph{funnel}, no la eficiencia de conversión.
\end{example}


%% ============================================================
\section{Plataformas de publicidad}

Las métricas de plataformas publicitarias viven dentro de los mecanismos de subasta de los
canales de pago. Diagnostican la eficiencia y cobertura de la presencia de la empresa dentro
de una plataforma específica, e informan la estrategia de puja más que la estrategia de
\emph{marketing} global.

\subsection{eCPM --- Effective Cost Per Thousand Impressions}

\textbf{Fórmula.}
\[
  \mathrm{eCPM} = \frac{\text{Total cost}}{\text{Total impressions}} \times 1000
\]

\textbf{Qué significa cada término.}
\begin{description}
  \item[Total cost] Gasto publicitario total, cualquiera que sea el tipo de compra utilizado
    --- costo por clic, costo por acción o CPM directo --- para la campaña o el período.
  \item[Total impressions] Todas las impresiones servidas bajo ese gasto, independientemente
    del tipo de compra.
\end{description}

\textbf{Qué mide.} El costo realizado total por mil impresiones en múltiples tipos de compra.
El eCPM normaliza diferentes mecanismos de subasta (compras CPC, CPA, CPM) en una única
cifra comparable de costo de alcance.

\textbf{Por qué importa.} Cuando un negocio ejecuta campañas bajo diferentes tipos de compra
--- algunas optimizadas para clics, otras para conversiones, otras para alcance --- el eCPM
proporciona una comparación consistente entre campañas de cuán eficientemente está comprando
audiencia cada una. También es la métrica de ingresos del editor: el lado de la oferta de la
subasta programática siempre cotiza en términos de eCPM.

\textbf{Referencia.} Usar el eCPM solo para comparar dentro del mismo tipo de posición y
segmento de audiencia. Las posiciones premium (bloques de portada, en el feed de redes
sociales) tienen eCPMs estructuralmente más altos que el display a nivel general de la red,
haciendo engañosas las comparaciones de eCPM entre posiciones.

\textbf{Trampa.} Comparar el eCPM en posiciones o audiencias muy diferentes trata a todas
las impresiones como equivalentes cuando no lo son. Un eCPM de \money{1.50} en display de
inventario residual y un eCPM de \money{60} en el feed ejecutivo de LinkedIn pueden ser
ambos apropiados para sus respectivos propósitos; la comparación no dice nada útil sobre la
eficiencia relativa en términos de resultado.

\begin{example}[eCPM en la SaaS de referencia]
Una campaña de Google de compra mixta: gasto total \money{30000}; \num{2000000} impresiones
entregadas entre Search (\num{400000}) y Display (\num{1600000}).
\[
  \mathrm{eCPM} = \frac{\money{30000}}{2000000} \times 1000 = \money{15}.
\]
Solo las impresiones de Search: gasto \money{20000} / \num{400000} impresiones $\times$ 1000
$= \money{50}$ eCPM. Solo Display: \money{10000} / \num{1600000} $\times$ 1000
$= \money{6.25}$ eCPM. El combinado de \money{15} oculta la diferencia estructural ---
comparar el eCPM de Search solo contra otras campañas de Search, y el de Display solo
contra Display.
\end{example}


\subsection{Impression Share --- Cuota de impresiones}

\textbf{Fórmula.}
\[
  \text{Impression Share} = \frac{\text{Impressions won}}{\text{Impressions eligible}}
\]

\textbf{Qué significa cada término.}
\begin{description}
  \item[Impressions won] El número de subastas en las que el anuncio se mostró efectivamente
    a un usuario.
  \item[Impressions eligible] El número estimado de subastas en las que el anuncio era
    elegible para participar, basado en su configuración de segmentación, presupuesto y
    señales de calidad.
\end{description}

\textbf{Qué mide.} Qué fracción de las impresiones elegibles disponibles está capturando
realmente la campaña. Una cuota de impresiones por debajo del \qty{100}{\percent} significa
que la campaña está perdiendo subastas para las que está calificada --- bien porque el
presupuesto se agotó (IS limitada por presupuesto) o porque la calidad o el rango del anuncio
fue demasiado bajo (IS limitada por rango).

\textbf{Por qué importa.} La cuota de impresiones diagnostica el cuello de botella en la
cobertura de pago. Si la IS es baja por presupuesto, añadir presupuesto aumentará directamente
la cobertura. Si la IS es baja por rango (calidad del anuncio o puja), aumentar el presupuesto
no sirve de nada; se requiere mejora del Índice de Calidad o aumento de la puja. Dividir la
IS perdida en sus dos componentes (presupuesto vs rango) antes de cambiar la estrategia
previene el error frecuente de añadir presupuesto a una campaña limitada por rango.

\textbf{Referencia.} El objetivo saludable depende de la intensidad competitiva. Perseguir
una IS del \qty{100}{\percent} en palabras clave amplias es una trampa de presupuesto;
una IS del \qty{80}{\percent}+ en los términos de mayor intención y mayor conversión es un
objetivo más defendible.

\textbf{Trampa.} Apuntar a una cuota de impresiones del \qty{100}{\percent} a cualquier
precio. En el margen, ganar cada impresión posible requiere superar las pujas de todos los
competidores en todos los términos, incluidos los de baja intención que convierten poco. Las
impresiones incrementales en el margen cuestan más y convierten menos que el promedio ---
rendimientos decrecientes clásicos de la maximización de IS.

\begin{example}[Cuota de impresiones en la SaaS de referencia]
Campaña de Google Search dirigida a «software de gestión de proyectos» y cinco términos
relacionados. Impresiones ganadas: \num{32000}; impresiones elegibles estimadas: \num{50000}.
\[
  \text{IS} = \frac{32000}{50000} = \qty{64}{\percent}.
\]
Desglose de IS perdida: \qty{20}{\percent} perdida por presupuesto, \qty{16}{\percent}
perdida por rango. Estrategia: aumentar el presupuesto para capturar primero el \qty{20}{\percent}
de IS perdida por presupuesto (inmediato, controlable) y ejecutar pruebas de texto de anuncio
para mejorar el Índice de Calidad para la parte perdida por rango (más lento, pero reduce el
CPC promedio).
\end{example}


\subsection{ACOS --- Advertising Cost of Sales}

\textbf{Fórmula.}
\[
  \mathrm{ACOS} = \frac{\text{Ad spend}}{\text{Attributed sales revenue}}
\]

\textbf{Qué significa cada término.}
\begin{description}
  \item[Ad spend] El costo publicitario del producto o campaña medida.
  \item[Attributed sales revenue] Ingresos por ventas atribuidos al anuncio, típicamente
    dentro de una ventana de atribución definida. El ACOS es la métrica principal en la
    plataforma publicitaria de Amazon; la ventana de atribución la define la plataforma.
\end{description}

\textbf{Qué mide.} El costo publicitario expresado como proporción de las ventas atribuidas
--- el inverso del ROAS. Un ACOS del \qty{25}{\percent} significa que se gastaron \money{0.25}
en publicidad por cada \money{1.00} de ventas.

\textbf{Por qué importa.} El ACOS es la métrica de eficiencia nativa de Amazon Sponsored
Products y plataformas de medios minoristas similares. Se usa para determinar si el gasto
publicitario de un producto está dentro del margen de umbral de rentabilidad. Como es el
inverso del ROAS, la interpretación es «menor es mejor» --- a diferencia del ROAS donde
mayor es mejor.

\textbf{Referencia.} El objetivo de ACOS $=$ margen bruto del producto. Si el margen del
producto es \qty{35}{\percent}, un ACOS objetivo del \qty{35}{\percent} significa que el
gasto publicitario iguala el beneficio bruto --- el producto alcanza el umbral de
rentabilidad a nivel de contribución. Un ACOS por debajo del objetivo de margen genera
beneficio bruto de la publicidad; por encima, destruye margen.

\textbf{Trampa.} Establecer objetivos de ACOS sin anclarlos al margen del producto. Un ACOS
del \qty{20}{\percent} en un producto con \qty{15}{\percent} de margen significa que cada
venta publicitada destruye dinero; el mismo ACOS del \qty{20}{\percent} en un producto con
\qty{50}{\percent} de margen es altamente rentable. El objetivo debe derivarse de la economía
del producto específico.

\begin{example}[ACOS para un caso ilustrativo de vendedor en Amazon]
Una \textit{marca de hardware de consumo} (ilustrativo, no SaaS): el producto se vende a
\money{89}; COGS \money{32}; margen bruto \qty{64}{\percent}. Gasto en Sponsored Products:
\money{4450}; ventas atribuidas: \money{22250}.
\[
  \mathrm{ACOS} = \frac{\money{4450}}{\money{22250}} = \qty{20}{\percent}.
\]
Dado que \qty{20}{\percent} $<$ margen del \qty{64}{\percent}, la campaña genera
$(\qty{64}{\percent} - \qty{20}{\percent}) = \qty{44}{\percent}$ de margen bruto sobre las
ventas publicitadas --- publicidad rentable. El ACOS de umbral de rentabilidad es
\qty{64}{\percent}; el negocio tiene margen significativo antes de que la campaña destruya
valor.
\end{example}


\subsection{iROAS --- Incremental Return on Ad Spend}

\textbf{Fórmula.}
\[
  \mathrm{iROAS} = \frac{\text{Incremental revenue caused by the ad}}{\text{Ad spend}}
\]

\textbf{Qué significa cada término.}
\begin{description}
  \item[Incremental revenue] Ingresos que \emph{no habrían ocurrido} sin la publicidad ---
    medidos mediante un experimento de retención (prueba geográfica, retención de usuarios o
    \emph{switchback}) que aísla la causalidad de la correlación.
  \item[Ad spend] El costo de la publicidad en la celda de prueba.
\end{description}

\textbf{Qué mide.} El retorno causal del gasto publicitario --- cuántos ingresos adicionales
causó genuinamente el anuncio, frente a los ingresos que habrían ocurrido de forma orgánica
o que fueron atribuidos por el modelo de último clic de la plataforma con independencia del
papel del anuncio.

\textbf{Por qué importa.} El iROAS es el estándar de oro para la medición publicitaria
porque responde la única pregunta que importa para la asignación de presupuesto: si gastamos
\money{1} más en este canal, ¿recupera el negocio más de \money{1}? El ROAS reportado por la
plataforma no puede responder esta pregunta porque atribuye conversiones orgánicas y asistidas
a los medios de pago. La metodología de medición se desarrolla completamente en
\cref{ch:attribution-and-measurement}.

\textbf{Referencia.} El iROAS debe superar 1 para que el gasto sea positivo en flujo de
caja. Debe superar $1/\text{margen bruto}$ para que el gasto sea positivo en beneficio bruto.
Para la SaaS de referencia con margen bruto aproximado del \qty{79}{\percent},
$\mathrm{iROAS}_{\min} = 1/0.79 \approx 1.27$.

\textbf{Trampa.} El iROAS requiere una retención o prueba geográfica --- no puede leerse en
los \emph{dashboards} de la plataforma. Sin un experimento, el iROAS no es observable. Las
afirmaciones de iROAS derivadas de modelos de atribución (incluso modelos sofisticados de
MMM) son salidas del modelo, no mediciones causales; conllevan incertidumbre del modelo que
debe divulgarse junto a la cifra.

\begin{example}[iROAS para la SaaS de referencia]
Una prueba de retención geográfica: los mercados de prueba reciben \money{40000} en gasto de
marca incremental durante seis semanas; los mercados de control reciben el gasto normal.
Ingresos en mercados de prueba: \money{1800000}; ingresos del mercado de control en
contrafactual (escalados): \money{1620000}. Ingresos incrementales: \money{180000}.
\[
  \mathrm{iROAS} = \frac{\money{180000}}{\money{40000}} = 4.5.
\]
ROAS reportado por la plataforma para la misma campaña: $7\times$ --- una sobreestimación del
\qty{56}{\percent}. El iROAS de 4.5 sigue superando el piso de margen bruto de 1.27,
confirmando que el gasto de marca crea valor; pero la decisión presupuestaria habría sido
distorsionada por la cifra reportada por la plataforma.
\end{example}


%% ============================================================
\section{Producto y engagement}

Las métricas de producto miden si los usuarios encuentran valor duradero en el producto, no
solo si lo visitan. Un motor de adquisición que llena un producto con fugas no es un motor
de crecimiento --- es una fábrica de \emph{churn}. Estas métricas diagnostican la sujeción
del producto sobre el comportamiento del usuario y la durabilidad de la relación con el
cliente.

\subsection{Stickiness --- Razón DAU/MAU}

\textbf{Fórmula.}
\[
  \text{Stickiness} = \frac{\mathrm{DAU}}{\mathrm{MAU}}
\]

\textbf{Qué significa cada término.}
\begin{description}
  \item[DAU] Usuarios Activos Diarios: el recuento de usuarios únicos que realizaron al menos
    una acción significativa en el producto en un día determinado. «Significativa» debe
    definirse (un inicio de sesión suele ser una señal demasiado débil; una interacción con
    una funcionalidad central es más defendible).
  \item[MAU] Usuarios Activos Mensuales: el recuento de usuarios únicos que realizaron al
    menos una acción significativa en el producto en cualquier momento del mes calendario.
\end{description}

\textbf{Qué mide.} La fracción de usuarios activos mensuales que regresan diariamente ---
un indicador directo de cuán habitual e integrado se ha vuelto el producto en los flujos de
trabajo diarios de los usuarios.

\textbf{Por qué importa.} La adherencia predice la retención. Un producto con
\qty{50}{\percent} de adherencia está genuinamente integrado en la rutina diaria; con
\qty{10}{\percent}, los usuarios visitan ocasionalmente y están a una mala experiencia de
sufrir \emph{churn}. Para SaaS B2B, la adherencia también impulsa los ingresos de expansión:
los equipos que usan un producto diariamente tienen mucha más probabilidad de añadir asientos
y hacer mejoras de plan que los usuarios ocasionales.

\textbf{Referencia.}
\begin{itemize}
  \item $>\qty{50}{\percent}$: élite; característico de productos de utilidad diaria (Slack,
    Gmail, Notion).
  \item $>\qty{20}{\percent}$: aceptable para herramientas B2B no diseñadas para uso diario
    (herramientas de planificación trimestral, de reporte anual se sitúan naturalmente más
    bajo).
  \item $<\qty{10}{\percent}$: engagement bajo; típicamente predice un \emph{churn} alto y
    poca expansión.
\end{itemize}

\textbf{Trampa.} Las cuentas de bots y las cuentas de usuario duplicadas inflan el DAU y el
MAU de forma independiente, pero la razón puede seguir pareciendo saludable incluso cuando
el engagement real de los usuarios cae. Validar siempre los recuentos de usuarios activos
mediante deduplicación de inicios de sesión y acciones antes de reportar la adherencia.

\begin{example}[Adherencia en la SaaS de referencia]
En el mes de medición: MAU $=$ \num{5500} (de \num{6700} clientes pagantes; algunos equipos
comparten inicios de sesión o usan el producto con poca frecuencia). Promedio de usuarios
activos diarios a lo largo del mes: \num{1650}.
\[
  \text{Stickiness} = \frac{1650}{5500} = \qty{30}{\percent}.
\]
Esta cifra supera el umbral aceptable B2B. Los \num{1200} clientes pagantes fuera del MAU se
señalan para una intervención proactiva del equipo de éxito del cliente --- las cuentas de
bajo uso son el indicador adelantado del \emph{churn}, no el retrasado.
\end{example}


\subsection{D30 Retention --- Tasa de retención al día 30}

\textbf{Fórmula.}
\[
  \text{D30 Retention} = \frac{\text{Users active at day 30}}{\text{Cohort size at day 0}}
\]

\textbf{Qué significa cada término.}
\begin{description}
  \item[Users active at day 30] Usuarios de la cohorte de adquisición que realizaron al
    menos una acción significativa en el producto el día 30 (o dentro de una ventana definida
    alrededor del día 30, p. ej., días 27--33).
  \item[Cohort size at day 0] El número de usuarios que se incorporaron (o activaron) en la
    cohorte. Una «cohorte» es un grupo definido por un evento de inicio compartido ---
    típicamente la semana de registro o la semana del primer inicio de sesión.
\end{description}

\textbf{Qué mide.} La fracción de una cohorte de nuevos usuarios que sigue interactuando con
el producto 30 días después de incorporarse. La retención D30 es la primera señal duradera
de si el producto ha creado valor duradero para los nuevos usuarios --- la retención en días
tempranos (D1, D7) suele estar inflada por la novedad.

\textbf{Por qué importa.} La retención D30 es el predictor temprano más sólido del LTV a
largo plazo del cliente. Una cohorte que retiene al \qty{40}{\percent} en el día 30 tendrá
una curva de retención a largo plazo muy diferente a la que retiene al \qty{15}{\percent}.
También diagnostica el \emph{funnel} de incorporación: una caída en D30 concentrada en las
primeras dos semanas señala un problema de incorporación; una caída en las semanas tres y
cuatro señala un problema de realización de valor (los usuarios llegaron al producto pero no
encontraron su valor central). El análisis completo de la curva de retención está en
\cref{ch:digital-funnel-and-crm}.

\textbf{Referencia.} Altamente dependiente de la categoría de producto. Móvil de consumo:
una retención D30 del \qty{20}{\percent} es un piso razonable. SaaS B2B con incorporación
activa: una retención D30 del \qty{60}{\percent}+ es alcanzable e importante para que
funcione la economía unitaria.

\textbf{Trampa.} Promediar la retención D30 entre fuentes de adquisición enmascara la
diferencia de calidad entre canales. Las cohortes de SEO orgánico pueden retener al
\qty{70}{\percent} mientras las cohortes de redes sociales de pago de una campaña de
audiencia amplia retienen al \qty{25}{\percent}; la cifra combinada del \qty{45}{\percent}
no describe con precisión a ninguna cohorte y envía al equipo de adquisición una señal
engañosa sobre la calidad del canal de pago.

\begin{example}[Retención D30 en la SaaS de referencia]
Una cohorte de registro de febrero de \num{800} usuarios de prueba (procedentes de
\num{10000} sesiones $\to$ \num{800} pruebas). Activos en el día 30: \num{520}.
\[
  \text{D30 Retention} = \frac{520}{800} = \qty{65}{\percent}.
\]
Dentro de esta cohorte, los registros de Google orgánico (600): retención en el día 30 del
\qty{74}{\percent}. Registros de LinkedIn de pago (200): retención en el día 30 del
\qty{41}{\percent}. La cifra combinada del \qty{65}{\percent} oculta una diferencia de
calidad significativa que debería informar cómo se asigna el presupuesto mensual de
adquisición de \money{120000}.
\end{example}


\subsection{NPS --- Net Promoter Score}

\textbf{Fórmula.}
\[
  \mathrm{NPS} = \%\,\text{Promoters} - \%\,\text{Detractors}
\]
donde los Promotores son los encuestados que puntúan 9--10 en la pregunta de «recomendaría»,
y los Detractores puntúan 0--6.

\textbf{Qué significa cada término.}
\begin{description}
  \item[\% Promoters] Proporción de encuestados que otorgan una puntuación de 9 o 10 sobre
    10 en la pregunta: «¿Con qué probabilidad recomendaría [producto] a un colega?»
  \item[\% Detractors] Proporción de encuestados que otorgan una puntuación de 0--6
    (inclusive). Las puntuaciones de 7--8 son «Pasivos» y quedan excluidas del cálculo.
\end{description}

\textbf{Qué mide.} Un indicador proxy grueso y comprimido de la satisfacción del cliente y
la propensión al boca a boca. El NPS va de $-100$ (todos Detractores) a $+100$ (todos
Promotores).

\textbf{Por qué importa.} El NPS se usa ampliamente como indicador de primer nivel del
sentimiento del cliente y para la comparación con pares del sector. Es predictivo del
\emph{churn} a nivel de cohorte --- las cohortes de Detractores sufren \emph{churn} a tasas
materialmente más altas que las cohortes de Promotores --- y captura el componente de boca a
boca del modelo de crecimiento.

\textbf{Referencia.} Depende del sector. SaaS: por encima de 30 es positivo, por encima de
50 es sólido. Por debajo de 0 es una señal de alerta que requiere investigación cualitativa
inmediata.

\textbf{Trampa.} Leer una encuesta de 11 puntos como un número preciso. Un cambio en el NPS
de 42 a 47 en una muestra de 300 encuestados es estadísticamente indistinguible del ruido
--- el intervalo de confianza típicamente abarca 10--15 puntos. La dirección de la tendencia
a lo largo de los trimestres importa más que la precisión en un momento dado. Además: el
sesgo de respuesta a la encuesta representa sistemáticamente en exceso a los clientes muy
satisfechos y muy insatisfechos, perdiendo al grupo ambivalente del medio.

\begin{example}[NPS en la SaaS de referencia]
Encuesta NPS trimestral enviada a \num{6700} clientes; \num{1340} respuestas (tasa de
respuesta del \qty{20}{\percent}). Promotores (9--10): 670 (\qty{50}{\percent}). Pasivos
(7--8): 469 (\qty{35}{\percent}). Detractores (0--6): 201 (\qty{15}{\percent}).
\[
  \mathrm{NPS} = 50 - 15 = 35.
\]
Las 201 cuentas Detractoras se cruzan con los datos de uso: el \qty{78}{\percent} de ellas
tiene una adherencia inferior al \qty{10}{\percent}, confirmando el bajo engagement como
principal impulsor de la insatisfacción y priorizando la intervención del equipo de éxito
del cliente.
\end{example}


\subsection{Activation Rate --- Tasa de activación}

\textbf{Fórmula.}
\[
  \text{Activation Rate} = \frac{\text{Users who reach the ``aha moment''}}{\text{Total signups}}
\]

\textbf{Qué significa cada término.}
\begin{description}
  \item[«Aha moment»] La acción o hito específico dentro del producto que se correlaciona
    más fuertemente con la retención a largo plazo. Para una SaaS de gestión de proyectos,
    podría ser «creó el primer proyecto con 3+ colaboradores». Para CRM: «registró el primer
    acuerdo y programó un seguimiento». El momento «aha» se define empíricamente a partir
    del análisis de retención de cohortes, no por suposición.
  \item[Total signups] Todos los usuarios que se registraron durante el período de medición,
    incluidas pruebas y cuentas gratuitas.
\end{description}

\textbf{Qué mide.} La fracción de nuevos usuarios que alcanzan con éxito el momento de valor
central del producto durante la incorporación. La tasa de activación es la bisagra entre
adquisición y retención: los usuarios que se activan retienen a 3--5$\times$ la tasa de los
que no lo hacen.

\textbf{Por qué importa.} Mejorar la tasa de activación multiplica el valor de cada dólar
de adquisición: si 30 de cada 100 registros de prueba se activan, duplicar la activación a
60 duplica efectivamente el rendimiento de la adquisición sin gastar un dólar adicional en
medios. Es la intervención de incorporación de mayor apalancamiento disponible. El \emph{funnel}
de activación se conecta directamente con \cref{ch:digital-funnel-and-crm}.

\textbf{Referencia.} Las tasas de activación sólidas varían según la complejidad del
producto. Herramientas de consumo simples: \qty{60}{\percent}+ es alcanzable. Plataformas
B2B complejas: \qty{30}{\percent}--\qty{50}{\percent} es saludable; por debajo del
\qty{20}{\percent} típicamente señala un problema de incorporación suficientemente grave como
para anular la inversión en adquisición.

\textbf{Trampa.} Definir «activado» de manera demasiado laxa infla la métrica sin predecir
la retención. Si la activación se define como «inició sesión al menos una vez», casi todos
los registros «se activan» pero la métrica no predice nada. La definición correcta es el
evento que diferencia de forma máxima las curvas de retención a largo plazo de los usuarios
activados frente a los no activados.

\begin{example}[Tasa de activación en la SaaS de referencia]
El equipo de producto define la activación como: «creó un espacio de trabajo, añadió al
menos dos integraciones e invitó a un colaborador dentro de los primeros 7 días». De
\num{800} registros de prueba este mes, \num{296} alcanzaron este hito.
\[
  \text{Activation Rate} = \frac{296}{800} = \qty{37}{\percent}.
\]
El análisis de retención muestra que los usuarios activados retienen al \qty{82}{\percent}
en el día 30; los no activados al \qty{19}{\percent}. La tasa de activación del 37\% significa
que 504 usuarios de prueba están en la trayectoria de baja retención. Una secuencia de emails
de incorporación dirigida con un \qty{30}{\percent} de mejora en la activación añadiría
aproximadamente 151 usuarios activados por mes, elevando las conversiones a pago esperadas
de 200 a 238 sin gasto adicional en adquisición --- equivalente a un \qty{19}{\percent} de
incremento en el ROI efectivo de \emph{marketing}.
\end{example}


%% ============================================================
\section{Métricas de vanidad --- Lo que ocultan}

Las métricas de vanidad se mueven con el esfuerzo, no con los resultados. Cada una en esta
sección distorsiona activamente la toma de decisiones en lugar de ser simplemente imprecisa:
puede aumentar mientras el negocio se deteriora, y optimizar para ella acelera el
deterioro. Cada métrica que sigue lleva una \textbf{admonition} en lugar de un ejemplo
trabajado, más un indicador de la métrica de resultado que debería reemplazarla.

\subsection{Raw Impressions --- Impresiones brutas}

\textbf{Fórmula.}
\[
  \text{Raw Impressions} = \text{total ad views (served impressions)}
\]

\textbf{Qué significa cada término.}
\begin{description}
  \item[Served impressions] El número de veces que un anuncio se entregó técnicamente a un
    dispositivo. Esto incluye impresiones nunca vistas realmente por un humano (debajo del
    pliegue, reproducción automática sin verificación de visualizabilidad, tráfico no válido).
\end{description}

\textbf{Qué mide.} Gasto, no alcance. El volumen de impresiones brutas escala mecánicamente
con el presupuesto; no indica que un humano real viera el anuncio, que fuera relevante o que
produjera algún efecto.

\textbf{Por qué engaña.} Una campaña puede multiplicar el recuento de impresiones aumentando
el presupuesto, ampliando la segmentación o comprando inventario residual barato --- ninguno
de lo cual mejora los resultados del negocio. Usar \emph{alcance único con visualizabilidad}
(el número de humanos distintos que vieron una impresión verificada como visible) y
\emph{frecuencia} (impresiones por usuario único) en su lugar.

\begin{admonition}{Advertencia}
Las impresiones brutas pueden aumentar un \qty{300}{\percent} mientras el alcance efectivo
permanece plano --- simplemente comprando la misma audiencia tres veces en lugar de llegar
a tres veces más personas. Una campaña que reporta «500 millones de impresiones» podría
representar 50 millones de personas viendo el anuncio 10 veces cada una (frecuencia excesiva,
probablemente con rendimientos decrecientes a partir de la impresión 4) o 500 millones de
personas distintas viéndolo una vez (alcance amplio y de contacto ligero). El recuento de
impresiones es idéntico; la implicación para el negocio es opuesta.
Usar el alcance y la frecuencia como la unidad de reporte combinada, nunca las impresiones
solas.
\end{admonition}


\subsection{Email Open Rate --- Tasa de apertura de correo electrónico}

\textbf{Fórmula.}
\[
  \text{Open Rate} = \frac{\text{Opens}}{\text{Emails delivered}}
\]

\textbf{Qué significa cada término.}
\begin{description}
  \item[Opens] Tradicionalmente rastreado por una imagen de seguimiento de 1 píxel que se
    carga cuando el email se muestra. Desde que Apple lanzó la Protección de Privacidad del
    Correo (MPP) en septiembre de 2021, Apple Mail precarga este píxel en nombre del usuario
    independientemente de si abrió el email.
  \item[Emails delivered] Enviados menos los rebotes duros.
\end{description}

\textbf{Qué mide.} Casi nada desde la MPP de Apple. Para listas con uso significativo de
Apple Mail (\qty{50}{\percent}+), las tasas de apertura reportadas son en gran medida función
de cuántos suscriptores usan Apple Mail, no de si realmente leyeron el email.

\textbf{Por qué engaña.} Las tasas de apertura reportadas en la mayoría de los ESP se han
inflado sistemáticamente desde septiembre de 2021 por los proxies de privacidad de Apple.
Una lista que muestra una tasa de apertura del \qty{48}{\percent} puede tener una tasa de
apertura humana real del \qty{20}{\percent}. Optimizar las líneas de asunto para la tasa de
apertura usando estos datos optimiza para lo que el proxy de Apple recupera, no para el
comportamiento humano. Usar la \emph{tasa de clics} (clics por email entregado) y la
\emph{tasa de respuesta} como métricas de engagement primarias; para emails de respuesta
directa, usar \emph{ingreso por email enviado}.

\begin{admonition}{Advertencia}
Una SaaS B2B medía el rendimiento del email por tasa de apertura y celebró un aumento del
\qty{32}{\percent} al \qty{47}{\percent} tras una prueba de línea de asunto. Las tasas de
clics eran planas. El ingreso por email enviado era plano. La mejora en la tasa de apertura
era completamente atribuible a un cambio en la composición de la lista hacia usuarios de
Apple Mail --- el proxy infló el denominador a una tasa diferente. Las decisiones tomadas
con estos datos reasignaron el presupuesto hacia el email y alejándolo de canales con datos
medibles de clic y conversión, dañando el ROI general de \emph{marketing}. Reportar las
tasas de apertura únicamente como un artefacto histórico; nunca tomar decisiones de
asignación de recursos a partir de ellas.
\end{admonition}


\subsection{Last-Click ROAS --- ROAS de último clic}

\textbf{Fórmula.}
\[
  \text{Last-Click ROAS} = \frac{\text{Revenue attributed to last click}}{\text{Ad spend}}
\]

\textbf{Qué significa cada término.}
\begin{description}
  \item[Revenue attributed to last click] Ingresos de las conversiones en las que el canal
    reportado recibió el último clic antes de la compra. Todos los canales anteriores
    (notoriedad, consideración, asistencias de retargeting) reciben cero crédito.
  \item[Ad spend] El costo de la campaña que se mide.
\end{description}

\textbf{Qué mide.} Quién cerró la venta, no quién la causó. El ROAS de último clic
identifica el canal que resultó estar en contacto en el momento de la conversión, que en un
recorrido de múltiples puntos de contacto suele ser una campaña de búsqueda de marca o
retargeting que interceptó una demanda generada por un punto de contacto anterior.

\textbf{Por qué engaña.} La atribución de último clic sobreacredita sistemáticamente al
último punto de contacto (búsqueda de marca, directo, retargeting) y subacredita a los
canales que crearon la notoriedad e intención iniciales (prospección en redes sociales,
contenido, PR). El presupuesto optimizado en ROAS de último clic traslada el gasto hacia
los canales de \emph{captura} de demanda y desfinancia los canales de \emph{creación} de
demanda, agotando eventualmente la parte alta del \emph{funnel} y causando que el propio ROAS
de último clic decaiga a medida que se erosiona la base de notoriedad subyacente. Usar
\emph{iROAS de experimentos de incrementalidad} como alternativa causal
(\cref{ch:attribution-and-measurement}).

\begin{admonition}{Advertencia}
Una conocida marca de comercio electrónico optimizó exclusivamente en ROAS de último clic
durante 18 meses. El ROAS de búsqueda de marca aumentó de $6\times$ a $11\times$ (estaba
capturando la demanda creada por un capital de marca construido durante años).
Simultáneamente, el gasto en la parte alta del \emph{funnel} se redujo del \qty{35}{\percent}
al \qty{12}{\percent} del presupuesto. El volumen de adquisición de nuevos clientes cayó un
\qty{40}{\percent} durante el período; la base de clientes existente envejeció y los ingresos
declinaron dos años después. La cifra de ROAS de último clic era la más alta que había sido
jamás el día en que el tráfico orgánico de búsqueda de la marca comenzó su declive
estructural. Las pruebas de iROAS, realizadas demasiado tarde, mostraron que los canales de
la parte alta del \emph{funnel} tenían un iROAS de 3.8 --- altamente rentable y privados de
presupuesto hasta cero.
\end{admonition}


\subsection{Unqualified Leads --- Prospectos no calificados}

\textbf{Fórmula.}
\[
  \text{Unqualified Leads} = \text{form fills or contact requests}
\]

\textbf{Qué significa cada término.}
\begin{description}
  \item[Form fills] Cualquier envío de un formulario web, descarga de contenido restringido,
    clic en solicitud de demo o solicitud de contacto, independientemente de la intención de
    compra, el encaje con la empresa o la autoridad de compra.
\end{description}

\textbf{Qué mide.} Interés, no demanda. Un prospecto no calificado ha expresado alguna forma
de interés pero no ha sido evaluado contra los criterios que hacen a un prospecto un comprador
viable (tamaño de empresa, presupuesto, autoridad, necesidad, timing --- BANT o un marco
equivalente).

\textbf{Por qué engaña.} El volumen de prospectos no calificados escala con el gasto en la
parte alta del \emph{funnel} y el volumen de contenido restringido, ambos de los cuales
pueden aumentarse mecánicamente sin ninguna mejora en la calidad del \emph{pipeline}. Un
equipo de contenido que publica una descarga viral de «Calculadora de ROI Gratis» puede
generar 5,000 formularios en un mes con cero intención de compra. Reportar el volumen de
prospectos a la dirección crea presión para generar más prospectos, no mejores, lo que
resulta en equipos de ventas sepultados en contactos de baja calidad. Devolver \emph{el
valor del pipeline calificado} (oportunidades calificadas por ventas con ingresos adjuntos)
al equipo de \emph{marketing} en su lugar, y atribuir ingresos contra gasto en lugar de
prospectos contra gasto. Véase \cref{ch:paid-acquisition-platforms} para la mecánica
completa de conversión de CPL a CPA.

\begin{admonition}{Advertencia}
Un equipo de \emph{marketing} de SaaS B2B era medido por el volumen mensual de prospectos.
Para alcanzar 2,000 prospectos por mes, el equipo compró una lista de emails de terceros y
promocionó un libro blanco genérico. Los prospectos llegaron a 2,500 --- un récord. Prospectos
aceptados por ventas (SAL): 48. Ganados y cerrados: 3. Costo por cliente cerrado desde esta
campaña: más de \money{15000}, frente al CAC combinado de \money{600}. El número de
prospectos subió; el negocio fue perjudicado. Medir \emph{el pipeline calificado generado}
y los \emph{ingresos influenciados} en su lugar --- métricas en las que los equipos de ventas
y \emph{marketing} deben ponerse de acuerdo antes del lanzamiento de la campaña.
\end{admonition}


\subsection{Follower Count --- Número de seguidores}

\textbf{Fórmula.}
\[
  \text{Follower Count} = \text{total followers on a social platform}
\]

\textbf{Qué significa cada término.}
\begin{description}
  \item[Followers] Cuentas de usuario que han hecho clic en «Seguir» en el perfil social de
    la marca. En la mayoría de las plataformas, la mayoría de los \emph{followers} ven menos
    del \qty{5}{\percent} de las publicaciones orgánicas debido a los límites de distribución
    algorítmica.
\end{description}

\textbf{Qué mide.} Tamaño de audiencia en una plataforma, no valor de audiencia. El número
de seguidores mide una acumulación histórica de clics en «Seguir» --- no mide el engagement,
el alcance del contenido actual ni ningún vínculo con los ingresos.

\textbf{Por qué engaña.} El número de seguidores puede crecer mediante tácticas de
seguir/dejar de seguir, adquisición pagada de seguidores, contenido viral fuera de tema y
campañas de sorteos --- todos los cuales añaden seguidores sin ninguna relación comercial
con la marca. Una marca con un millón de seguidores que alcanza al \qty{0.5}{\percent}
orgánicamente (5,000 personas) puede ser menos efectiva que un competidor con 50,000
seguidores altamente comprometidos que alcanza al \qty{20}{\percent} (10,000 personas).
Rastrear \emph{el alcance con engagement por publicación}, la \emph{tasa de conversión del
tráfico de redes sociales} y los \emph{ingresos atribuibles al contenido de redes sociales}
en su lugar.

\begin{admonition}{Advertencia}
Una empresa de software B2B hizo crecer su número de seguidores de LinkedIn de 8,000 a
42,000 en seis meses mediante una serie de contenido viral de «memes de la industria».
Tráfico web desde LinkedIn: plano. Solicitudes de demo desde LinkedIn: sin cambios. Los
seguidores adquiridos eran en su mayoría colaboradores individuales de industrias adyacentes
que nunca comprarían el producto. La estrategia de contenido fue posteriormente reestructurada
en torno a publicaciones de liderazgo de pensamiento dirigidas al perfil real del comprador
(VP de Operaciones, Jefe de Gabinete), reduciendo la tasa de crecimiento de seguidores pero
aumentando la atribución de solicitudes de demo desde LinkedIn en un \qty{220}{\percent}.
\end{admonition}


\subsection{App Downloads --- Descargas de aplicación}

\textbf{Fórmula.}
\[
  \text{App Downloads} = \text{total installs from app store}
\]

\textbf{Qué significa cada término.}
\begin{description}
  \item[Downloads / Installs] El recuento de veces que la aplicación fue instalada en un
    dispositivo. Esto incluye instalaciones abandonadas de inmediato, instalaciones en
    dispositivos que nunca se abren y reinstalaciones por el mismo usuario tras eliminar la
    aplicación.
\end{description}

\textbf{Qué mide.} Prueba, no adopción. Una descarga es el inicio del recorrido del cliente,
no un punto en la curva de retención.

\textbf{Por qué engaña.} Los recuentos de descargas en la tienda de aplicaciones pueden
inflarse mediante trucos de optimización de la tienda, campañas de instalación de audiencia
amplia e instalaciones incentivadas («descarga para ganar un premio»). Una aplicación de
consumo con 10 millones de descargas y una retención D30 del \qty{5}{\percent} tiene 500,000
usuarios; la misma aplicación con 1 millón de descargas y una retención D30 del
\qty{40}{\percent} tiene 400,000 usuarios. La primera parece 10$\times$ mayor; los negocios
son comparables en usuarios activos. Reportar \emph{la tasa de activación}, \emph{la
retención en el día 30} y \emph{los usuarios activos mensuales} junto a los recuentos de
descargas; nunca usar las descargas como métrica de éxito de forma aislada.

\begin{admonition}{Advertencia}
Una \emph{startup} de juegos para móvil celebró 5 millones de descargas en sus primeros tres
meses, atrayendo una extensión de \emph{seed} a una valoración de \money{15000000} anclada
en la velocidad de descargas. La retención D30 era del \qty{3}{\percent}; MAU al mes cuatro:
\num{48000}. Al mes ocho, el DAU había caído por debajo de \num{10000} y la velocidad de
nuevas instalaciones se había estancado a medida que la campaña de instalación de audiencia
amplia agotó el mercado direccionable. La métrica de vanidad impulsó una decisión de
financiación basada en premisas falsas; la métrica de resultado (retención D30 $\times$ tamaño
de cohorte) estaba disponible desde el primer día y predijo la trayectoria con precisión.
\end{admonition}
